\chapter{Od výpočtu k nadvládě: pojmové základy politické analýzy umělé inteligence}
\section{Umělá inteligence jako předmět filozoficko-politické analýzy.}
Historie ideje umělého rozumu je podstatně starší, než se obvykle předpokládá. 
Nezačíná Dartmouthskou konferencí roku 1956 a není výlučným produktem technologického pokroku dvacátého století.
 Naopak, za touto ideou stojí mnohasetletá filozofická tradice, v níž každá epocha po svém odpovídala na tutéž otázku: 
 co je rozum a lze jej reprodukovat?

Starověká řecká kultura znala mechanické divy: Hérón Alexandrijský popisoval složité divadelní a chrámové automaty,
 mytologická tradice přisuzovala Héfaistovi schopnost vytvářet mechanické bytosti.
 Zásadně důležité je však to, že řecké chápání světa nepřipouštělo možnost autonomního umělého rozumu.
  Rozum - nous - byl chápán jako univerzální, kosmický princip, v zásadě nedostupný pro reprodukci v jednotlivém artefaktu. 
 Hérakleitos přímo tvrdil, že mezi lidskou a božskou přirozeností rozumu leží nepřekročitelná propast. 
 V tomto souřadnicovém systému nebyly antické mechanismy nikdy pojímány jako myslící bytosti 
 — byly pouhými dovedně provedenými imitacemi pohybu, nikoli rozumu. 
 Toto omezení nebylo technické, nýbrž ontologické: plynulo ze samotného chápání toho, čím je intelekt.

\subsection{Středověk: první projekt mechanizace myšlení.}

Zásadně odlišnou intelektuální možnost otevřelo křesťanský středověk. 
Křesťanská teologie \textbf{trvala} na tom, že člověk byl stvořen k obrazu a podobě Boží a nese v sobě jiskru rozumu blízkého rozumu božskému.
 Z toho plynul nečekaný důsledek: je-li rozum člověku darován Bohem, pak reprodukce rozumu v artefaktu stvořeném člověkem přestává být rouháním a stává se pokračováním tvůrčího aktu.

Právě v tomto kontextu je třeba chápat projekt Raimunda Lulla (cca 1232–1316). 
Jeho Ars Magna — kombinatorický systém založený na mechanickém generování pravdivých soudů prostřednictvím manipulace se symbolickými prvky — byl historicky prvním systematickým projektem formalizace myšlení. 
Ačkoli byl svými cíli teologický, metodologicky byl revoluční: poprvé byl vysloven předpoklad, že rozum není mystická schopnost dostupná pouze vyvoleným, nýbrž vypočitatelná operace nad symboly.
 V tomto smyslu je Lull přímým předchůdcem Leibnizovy characteristica universalis, Boolovy algebry logiky a v konečném důsledku i Turingova výpočetního modelu.

Paralelně se ve středověké kultuře rozvíjela kabalistická tradice golema — umělé bytosti oživené prostřednictvím posvátných slov.
 Strukturální podobnost s Lullovým projektem je zde zásadní: v obou případech je rozum reprodukován správnou manipulací se symboly nadanými zvláštní mocí.
  Dlouho před kybernetikou středověké myšlení intuitivně nahmátlo tezi, která se stane ústřední pro teorii informace: rozum je struktura, informace je forma.

\subsection{Od mechanicismu ke kybernetice: tři konceptuální posuny.}

Přechod od středověkého mýtu k vědeckému projektu UI se uskutečňoval prostřednictvím několika konceptuálních posunů, z nichž každý je svým způsobem zásadní.

První posun — karteziánský mechanicismus.
 Descartes rozdělil svět na substanci rozlehlou (res extensa) a myslící (res cogitans). 
Tělo je mechanismus podřízený zákonům fyziky; myšlení je něco v zásadě odlišného. 
Toto rozdělení zrodilo dualitu, která bude pronásledovat veškerý následující diskurz o UI: na jedné straně je stroj v zásadě oddělen od myšlení; na straně druhé mechanizované tělo navozuje otázku, zda i myšlení není v jistém smyslu mechanickým procesem. 
Leibniz se vydal právě tímto směrem: jeho characteristica universalis předpokládala, že všechny pravdy lze vyjádřit v univerzálním jazyce a ověřit mechanicky prostřednictvím calculus ratiocinator. Myšlení bylo redukováno na formální operaci.

Druhý posun — průmyslová revoluce a mechanicismus 19. století. 
Vytvoření stále složitějších výpočetních zařízení, jako byl Babbageův analytický stroj či Jacquardovy děrné štítky, učinilo ideu myslícího stroje technicky hmatatelnou. Zároveň materialistická filozofie stále naléhavěji kladla otázku, zda mozek sám není druhem výpočetního zařízení. Tato otázka přestala být čistě spekulativní a získala inženýrský rozměr.

Třetí a rozhodující posun — kybernetická revoluce poloviny 20. století.
 Norbert Wiener v Kybernetice \parencite{Wiener48} navrhl univerzální jazyk pro popis řízení a komunikace v systémech jakékoli povahy, ať živých či mechanických. Klíčové koncepty — informace, zpětná vazba, řízení — umožnily nahlížet na rozum jako na informační proces. Na tomto základě vybudoval Turing svou argumentaci: je-li myšlení výpočtem, pak jakýkoli systém schopný dostatečně složitých výpočtů je schopen i myšlení \parencite{Turing50}. Dartmouthská konference roku 1956 tuto tezi zakotvila v podobě výzkumného programu.

Právě zde však, na vrcholu technologického optimismu, vyvstala zásadní filozofická kontroverze.
 Hubert Dreyfus v díle What Computers Can't Do \parencite{Dreyfus92} ukázal, že symbolická UI přehlíží fundamentální rys lidského poznání — jeho vtělenost: intelekt je neoddělitelný od tělesné zkušenosti, praktického vědění a situovanosti ve světě, jež nepodléhají úplné formalizaci.

 John Searle ve svém myšlenkovém experimentu \uv{Čínský pokoj}  \parencite{Searle80} k tomu přidal rozlišení, jež se stalo klasickým: výpočetní systémy manipulují symboly podle formálních pravidel, ale nerozumí jim — mezi syntaxí a sémantikou leží propast, kterou žádná sebesložitější program překonat nedokáže. 
 Systém může produkovat syntakticky správná a pragmaticky přesvědčivá tvrzení, aniž by rozuměl jejich obsahu.

 Tato námitka neztratila aktuálnost s příchodem moderních systémů hlubokého učení. 
 Systémy, které nerozumí, ale přesvědčivě napodobují porozumění, jsou dnes zabudovány do mechanismů správy, moderace veřejného diskurzu a formování veřejného mínění. Jak poukazuje Floridi, právě to činí vypracování normativních rámců pro UI naléhavým úkolem: technologie zbavená sémantického porozumění, avšak nadaná sociální mocí, vyžaduje vnější etická a politická omezení, která sama produkovat není schopna \parencite{Floridi23}.

\section{Transformace pojmu \uv{myšlení}: od lidského k výpočetnímu.}
V antické a středověké tradici bylo myšlení chápáno jako něco v zásadě nesvoditelného na logický úsudek.
Aristotelův koncept fronésis \uv{praktická moudrost} označoval schopnost mravního úsudku v konkrétní situaci, jež nelze redukovat na aplikaci obecných pravidel.
Nous jako kontemplace nejvyšších pravd a fronésis jako situované mravní rozlišování tvořily společně to, co řecká tradice rozuměla pod rozumem: schopnost celistvého úsudku zakořeněného v zkušenosti, mravnosti a vtělenosti.
Toto chápání rozumu v zásadě vylučovalo možnost jeho úplné formalizace.

Karteziánský mechanicismus provedl první zúžení: myšlení začalo být ztotožňováno s jeho racionální, logickou složkou
 — tedy s tím, co potenciálně formalizaci podléhá.
Mravné, emocionální a estetické bylo buď vyloučeno, nebo redukováno na \uv{vášně}  podléhající řízení rozumem. 
Leibnizův projekt calculus ratiocinator učinil další krok: myšlení bylo ztotožněno s formálními operacemi nad symboly.
 Boolova logika a posléze matematická logika Frega a Russella tento proces dovršily na úrovni formálních systémů, v nichž byl rozum redukován na syntax.

Turing položil zásadní otázku: nemůžeme-li v dialogu rozlišit odpovědi stroje od odpovědí člověka, máme vůbec důvod popírat, že stroj myslí? \parencite{Turing50} 
Tato otázka byla revoluční — předpokládala redukci myšlení na pozorovatelné chování. Vnitřní zkušenost, intencionálnost, mravní zakořeněnost byly odsunuty jako nepozorovatelné, a tedy irelevantní pro operacionální definici inteligence. Searle ukázal principiální neudržitelnost této redukce: jeho myšlenkový experiment \uv{Čínský pokoj} demonstruje, že systém může produkovat syntakticky správné a pragmaticky přesvědčivé odpovědi, aniž by měl jakýkoli přístup ke smyslu manipulovaných symbolů — mezi syntaxí a sémantikou leží propast, kterou behaviorální test odhalit nedokáže \parencite{Searle80}. Dreyfus k tomu přidal argument vtělenosti: lidské poznání je neoddělitelné od tělesné zkušenosti, od praktického \uv{vědění-jak} na rozdíl od propozicionálního \uv{vědění-že}, a žádný formální systém tuto zakořeněnost reprodukovat nedokáže \parencite{Dreyfus92}.

Přesto se redukcionistická koncepce \uv{myšlení-jako-výpočtu} ukázala historicky vítěznou — je technicky produktivní a ekonomicky výhodná. 
Vládnoucí intuicí naší doby se stala intuice technokratická: rozum je to, co lze změřit, formalizovat a reprodukovat. 
Vše ostatní je buď nepodstatné, nebo bude formalizováno později. 
Floridi poukazuje, že právě tato redukce činí vypracování normativních rámců pro UI naléhavým úkolem: systém zbavený mravného a sémantického rozměru sám etická omezení neprodukuje — musejí být přinesena zvenčí \parencite{Floridi23}.

Důsledky tohoto konceptuálního vítězství dalece přesahují technické diskuse. 
Je-li myšlení výpočtem, může algoritmus přijímat rozhodnutí namísto člověka: o bonitě, o rizicích, o vině, o způsobilosti k funkci. 
Je-li inteligence měřena behaviorálním testem, pak systém, který testem projde, disponuje dostatečnou inteligencí pro řízení.
 Je-li mravný rozměr myšlení irelevantní, otázka etické odpovědnosti algoritmu jednoduše nevyvstane. 
 Olof-Ors a Smit ukazují, že právě toto konceptuální dědictví umožňuje antropomorfní design jako nástroj dozorového kapitalismu: systém, který přesvědčivě reprodukuje chování rozumné bytosti, je uživatelem jako taková vnímán, a celý aparát sociální důvěry, evolučně určený pro interakci s jinými lidmi, je tak nasazen do služeb korporace \parencite{OlofOrsSmit25}.

V historii od Lulla k moderním jazykovým modelům sledujeme postupné předefinování myšlení s cílem zahrnout do něj stroj.
 V každém stadiu bylo z tohoto pojmu něco podstatného vyloučeno: nejprve mravný úsudek, poté intencionálnost, poté vtělenost a nakonec samotné porozumění. Výsledkem je systém produkující statisticky pravděpodobné posloupnosti symbolů bez jakéhokoli přístupu k jejich smyslu — to, co by Searle nazval syntaxí bez sémantiky. To však v zásadě není to, co filozofická tradice myšlením nazývala. Na tomto rozlišení závisí, kdo nese odpovědnost za rozhodnutí přijímaná \uv{inteligenčními} systémy a čím zájmům tato rozhodnutí slouží.

\subsection{Od disciplinární moci k algoritmické kontrole.}

Michel Foucault v díle Discipline and Punish (1977) popsal přechod od svrchované moci, založené na právu na smrt, k moci disciplinární, založené na správě života prostřednictvím institucí — vězení, školy, nemocnice, kasárny. Gilles Deleuze v roce 1992 zaznamenal další posun: disciplinární společnosti ustupují společnostem kontroly, v nichž je moc vykonávána nikoli prostřednictvím uzavřených prostorů, nýbrž prostřednictvím kontinuálních, flexibilních a všudypřítomných modulací \parencite{Deleuze92}. Dnes můžeme pojmenovat konkrétní mechanismus této modulace — algoritmus.

Disciplinární moc v Foucaultově popisu fungovala prostřednictvím normalizace: jedinec je pozorován, měřen, porovnáván s normou a v případě potřeby korigován \parencite{Foucault77}. Benthamovo panoptikum bylo archetypem této logiky: k tomu, aby pozorovaný začal kontrolovat sám sebe, postačuje potencialita pozorování. Digitální panoptikum reprodukuje tuto strukturu jako flexibilní, inteligentní a prakticky nepostřehnutelná infrastruktura masového dohledu, umožňující korporacím sledovat jednání každého uživatele při zachování iluze naprosté svobody \parencite{SauraGarcia24}. Zásadní rozdíl oproti foucaultovskému panoptiku spočívá v tom, že jeho operátorem již není stát, nýbrž soukromá korporace, což radikálně mění otázku odpovědnosti a demokratické kontroly.

Zde se nabízí odkaz na klasický argument Langdona Winnera: artefakty mají politiku. 
Technická zařízení a systémy nejsou neutrálními nástroji, jež lze stejně dobře použít k různým účelům — strukturálně vtělují určité sociální vztahy a rozdělení moci \parencite{Winner86}. 
Algoritmické systémy jsou v tomto smyslu politicky nejnabitějšími artefakty současnosti: jejich architektura určuje, kdo co může vidět, kdo koho může nalézt, čí hlasy budou vyslyšeny a čí zůstanou ve stínu. 
Feenberg tuto teze rozvíjí v rámci kritické teorie technologie: technologie není neutrální silou stojící naproti společnosti, nýbrž samotným polem, na němž se odehrává zápas o to, jakou bude mít sociální realita podobu \parencite{Feenberg02}.

Akbari a Wood popisují toto jako strukturální transformaci autoritářství. 
Tradiční politické teorie se soustřeďovaly na stát jako hlavního nositele autoritářské moci; avšak dnes politici, zákonodárci ani voliči již nejsou jedinými determinanty politických rozhodnutí \parencite{AkbariWood25}. Platformové korporace se staly novou třídou autoritářských subjektů, jednajících prostřednictvím tří vzájemně provázaných dimenzí: autoritářských diskurzů, autoritářských praktik a autoritářské infrastruktury. Poslední dimenze je zvlášť důležitá: digitální infrastruktura — ať jde o algoritmus řazení zpravodajského kanálu, systém kreditního scoringu nebo Velký čínský firewall — představuje projektované materiální prostředí, v němž je moc vtělena strukturálně.

V tomto kontextu je vhodné připomenout Agambena. 
Jeho koncept výjimečného stavu \parencite{Agamben05} popisuje zónu, v níž jsou právní záruky pozastaveny mocí, jež si nárokuje jejich ochranu. 
Algoritmické systémy vytvářejí funkční analogy takového stavu: automatizovaná rozhodnutí o vízech, úvěrech, sociálních dávkách a moderaci obsahu jsou přijímána v zóně, kde jsou formální práva občana — na odůvodnění, na odvolání, na rovné zacházení — de facto pozastavena, neboť protistranou není instituce, nýbrž neprostupný technický systém. Shoshana Zuboff konceptualizuje ekonomickou logiku této transformace jako dozorový kapitalismus: lidská zkušenost je systematicky přeměňována v surovinu pro behaviorální predikce, prodávané na speciálně vytvořených trzích, což fundamentálně mění základy mocenských vztahů \parencite{Zuboff19}.
Algoritmus přitom tyto vztahy aktivně formuje prostřednictvím personalizovaného prostředí, v němž jsou přijímána rozhodnutí. 
Moc přestává potřebovat donucení: postačuje ovládat prostředí, v němž subjekt volí, a volba bude předvídatelná bez jakéhokoli explicitního zásahu.

\section{Depolitizace a mýtus algoritmické neutrality}

Ústředním politickým důsledkem algoritmizace moci je depolitizace: vytěsnění politického rozhodnutí — tj. rozhodnutí přijatého subjektem nesoucím za ně demokratickou odpovědnost — technickým výstupem algoritmu, jenž nemá ani tvář, ani adresu. Konceptuální jádro této transformace tvoří předpoklad, že složité sociální problémy lze přeměnit v jasně vymezené technické úlohy a řešit je algoritmicky \parencite{JanssenKuk16}. Koncepce \uv{chytrého města} je nejvýmluvnějším institucionálním příkladem: spočívá na předpokladu, že veškeré aspekty městského života lze měřit, kontrolovat a redukovat na technické problémy \parencite{Kitchin14}. Politická a sociální hlediska přitom ustupují do pozadí: jak řídící, tak řízení se ocitají podřízeni logice technologie, představované jako vnější vůči moci.

König a Wenzelburger ukazují, že to vytváří konkrétní strukturální důsledek pro demokratické instituce: algoritmické systémy nastolují podmínky, za nichž jsou politická rozhodnutí de facto přijímána na úrovni technické architektury — dříve, než se dostanou do prostoru veřejné diskuse \parencite{KoenigWenzelburger20}. Volba parametrů algoritmu kreditního scoringu, kritérií viditelnosti ve výsledcích vyhledávání, vah v systémech doporučení — to vše jsou v podstatě politická rozhodnutí, která určují rozdělení zdrojů, příležitostí a informací. Avšak jsou přijímána v uzavřených korporátních prostorách, nepodléhají veřejné diskusi a nejsou odpovědná voličům.

Coeckelbergh formuluje politicko-epistemologický důsledek tohoto procesu prostřednictvím konceptu epistemického agentství: demokracie vyžaduje reálnou schopnost občanů formovat odůvodněné soudy, kriticky hodnotit informace a odolávat manipulaci \parencite{Coeckelbergh23}. 
Algoritmické zprostředkování tyto schopnosti systematicky podkopává a jejich uplatňování stále více ztěžuje.
 Delegováním kognitivní práce na technický systém občan postupně ztrácí samotnou schopnost k nezávislému úsudku, a s ní i možnost plnohodnotně se účastnit demokratické politiky.

Saura García popisuje politický rozměr tohoto procesu prostřednictvím pojmu regulatorního zajatectví: za poslední desetiletí technologické korporace podstatně zvýšily výdaje na lobbying, rozšířily počet setkání svých vedoucích pracovníků se státními úředníky a vytvořily systémy \uv{otočných dveří} mezi regulatorními orgány a odvětvím \parencite{SauraGarcia24}.
Výsledkem je situace, kdy regulátoři de facto chrání zájmy těch, které mají regulovat. 
Floridi poukazuje, že to činí čistě technokratický přístup k regulaci UI strukturálně nedostatečným: regulace algoritmů vyžaduje politickou vůli, jež nemůže být delegována těm samým korporacím, jejichž systémy mají být regulovány \parencite{Floridi23}.

Williams rozvíjí blízkou argumentaci ve specifickém aspektu ekonomiky pozornosti: zmocnění se pozornosti platformami ničí individuální autonomii i infrastrukturu, na níž spočívá demokratická politika — tj. schopnost občanů soustředit svou pozornost na otázky vyžadující kolektivní diskusi, nikoli na to, co přináší platformám největší zisk \parencite{Williams18}. Když komerční infrastruktura určuje, na co občané hledí, na co reagují a o čem hovoří, demokratická politika ztrácí svůj materiální základ.

Paralelně probíhá to, co Olof-Ors a Smit nazývají obfuskací dozorového kapitalismu prostřednictvím antropomorfizace: přerámování vztahů s korporátními algoritmickými systémy jako sociálních vazeb, a nikoli komerčních transakcí, ztěžuje nejen technickou regulaci, ale i samotnou společenskou vůli ji vyžadovat \parencite{OlofOrsSmit25}.
Regulace takových systémů, jako je Replika, začíná být vnímána jako zásah do osobních vztahů. Depolitizace se tak uskutečňuje prostřednictvím emocionální architektury uživatelské zkušenosti.

\subsection{Algoritmická předpojatost a její politické důsledky}

Zdánlivá neutralita algoritmu patří k nejodolnějším mýtům technokratického diskurzu. 
Algoritmus se jeví jako objektivní nástroj zbavený lidských předsudků: zpracovává data a vydává výsledek, aniž by se řídil sympatiemi či ideologií. 
Právě tento obraz mu zajišťuje politickou legitimitu — rozhodnutí přijaté algoritmem bývá vnímáno jako spravedlivější než rozhodnutí přijaté člověkem. 
Tento obraz je však v zásadě mylný.

Cathy O'Neil ve Weapons of Math Destruction \parencite{ONeil16} formuluje: algoritmy dnes používané pro rozhodování o zaměstnávání, úvěrování, pojišťování, vzdělávání a trestním soudnictví představují \uv{zbraně matematického ničení}. Jedná se o modely neprůhledné ve své logice, rozsáhlé ve svých důsledcích a systematicky poškozující ty, kdo se již nacházejí v zranitelném postavení. O'Neil ukazuje, že algoritmy se dynamicky vyvíjejí spolu s daty, systémy a lidmi v rámci složitého sociotechnického prostředí a v tomto procesu vstřebávají předsudky obsažené v historických datech. Část systematických zkreslení je zděděna z databází odrážejících historické nerovnosti; jiná jsou nezáměrně vnášena vývojáři nebo vznikají v důsledku změn v informačním prostředí. Výsledkem je, že intelektuální analýza dat odhaluje diskriminační praktiky a zakořeněnou sociální nerovnost a posléze je reprodukuje, propůjčujíc jim zdání technické objektivity.

Safiya Umoja Noble v Algorithms of Oppression \parencite{Noble18} to dokládá na příkladu vyhledávačů: algoritmy Google systematicky reprodukují rasistické a sexistické stereotypy ve výsledcích vyhledávání, přičemž tento efekt nepředstavuje selhání systému — jde o strukturální výsledek komerční logiky, v níž jsou statistické korelace ve velkých datech přijímány za objektivní pravdu. Jelikož Google Search je fakticky monopolním zprostředkovatelem mezi občany a informacemi, taková systematická zaujatost dosahuje rozsahu politické infrastruktury.

Virginia Eubanks v Automating Inequality \parencite{Eubanks18} rozšiřuje tuto analýzu na oblast sociální politiky a ukazuje, jak automatizované systémy určování nároku na dávky, hodnocení rizika u dětí, přidělování bydlení a dalších sociálních statků systematicky diskriminují chudé rodiny v USA.
Eubanks potvrzuje: technologie prezentované jako nástroje zvyšování efektivnosti sociálních programů ve skutečnosti reprodukují a zesilují logiku sociální kontroly nad nejchudšími vrstvami obyvatelstva, sahající svými kořeny do 19. století. 
Zásadně důležité přitom je, že tyto systémy nejenže odrážejí stávající nerovnosti — nově definují, co znamená \uv{být v nouzi}, přeměňujíc politickou otázku spravedlivého přerozdělení v technický postup hodnocení rizika.

Ruha Benjamin v Race After Technology \parencite{Benjamin19} formuluje zobecňující koncept — \uv{nový Jimův zákon} (New Jim Code): technologie, jež jsou navenek neutrální, avšak ve skutečnosti reprodukují rasové hierarchie pod rouškou technické optimalizace. 
Stejně jako v předchozí epoše, kdy formálně neutrální zákony Jima Crowa systematicky diskriminovaly Afroameričany, mohou moderní algoritmické systémy rasu přímo nezmiňovat, avšak systematicky reprodukovat rasovou nerovnost prostřednictvím korelovaných proměnných: poštovního směrovacího čísla, vzdělání, úvěrové historie.

Saura García popisuje strukturální důsledek těchto procesů: velké asymetrie vědění mezi korporacemi a občany v kombinaci s monopolní kontrolou nad algoritmickými platformami umožňují korporacím provádět kampaně politického vlivu a manipulace prostřednictvím mechanismů mikrotargetingu, neurotargetingu, zkreslování informací a umělého formování veřejného mínění \parencite{SauraGarcia24}. Kauza Cambridge Analytica je v tomto ohledu nejznámějším příkladem. Akbari a Wood k tomu přidávají postkoloniální rozměr: autoritářské využití sledovacích technologií není výsadou oficiálně autoritářských režimů — dodavatelské řetězce sledovacích systémů lze vysledovat od USA, Francie, Německa a Izraele až po Nigérii, Maroko a Zambii \parencite{AkbariWood25}.

Algoritmická předpojatost zde nabývá geopolitického rozměru, v němž jsou sledovací technologie vyvinuté v demokratických zemích vyváženy do režimů, jež je používají proti vlastním občanům.

Zvláštní politickou naléhavost získává problém předpojatosti v oblastech bezprostředně spojených s rozdělováním životních příležitostí. Ve všech těchto případech přijímá algoritmus rozhodnutí dotýkající se konkrétního člověka, tento člověk však zpravidla neví o existenci algoritmu, nemá přístup k jeho logice a nemá k dispozici mechanismus odvolání. Peters ukazuje, že i v relativně příznivých podmínkách algoritmická politická zaujatost systematicky uniká technickým metodám jejího odhalování a nápravy, neboť je zabudována do samotné operacionalizace úlohy, kterou má algoritmus řešit \parencite{Peters22}. Floridi trvá na tom, že to je neslučitelné se základními požadavky demokratického vládnutí: odpovědnost předpokládá možnost vysledovat řetězec zodpovědnosti od rozhodnutí k osobě, která je přijala — algoritmus však tento řetězec přetrhuje \parencite{Floridi23}. Rozhodnutí delegovat rozhodnutí na algoritmus je samo o sobě politickou volbou — pouze skrytou.

Zde se vrací Agambenův koncept výjimečného stavu \parencite{Agamben05}: v zónách algoritmického rozhodování jsou práva subjektu formálně zachována, avšak fakticky pozastavena — odvolat se proti rozhodnutí \uv{černé skříňky} je nemožné ani právně, ani technicky.
\uv{Problém černé skříňky} lze sledovat v tom, že v složitých případech jsou algoritmy natolik komplexní, že ani jejich tvůrci nedokáží přesně vysvětlit, jak přesně fungují, přičemž podrobnosti jejich vývoje jsou chráněny jako obchodní tajemství. To znamená, že algoritmická předpojatost je systematicky neviditelná, a tedy i nezpochybnitelná.

Východisko z tohoto kruhu vyžaduje audit algoritmů či požadavky transparentnosti, jakož i politické přehodnocení samotné otázky, která rozhodnutí je vůbec přípustné delegovat na algoritmus. Jde o rozdělení moci: kdo rozhoduje o tom, jak je algoritmus sestaven, čím zájmům je optimalizován a kdo nese odpovědnost za jeho důsledky. Dokud tyto otázky zůstávají v prostoru technické expertízy, bude algoritmizace moci prohlubovat depolitizaci bez ohledu na to, jak dokonalé se ukáží být regulatorní nástroje.
\section{Nadvláda jako normativní jádro}
Dosavadní výklad popisoval, jak algoritmická moc funguje: jak se myšlení ztotožnilo s výpočtem a jak se politické rozhodnutí skrylo pod technickým výstupem. 
Zbývá však říci to podstatnější — v čem přesně spočívá její politické bezpráví.
 Odpověď, o niž se opírá zbytek práce, čerpá z neorepublikánské tradice politické filozofie, především z pojetí svobody jako nepřítomnosti nadvlády (non-domination) u Philipa Pettita a Franka Lovetta.
Republikánská tradice staví proti liberálnímu pojetí svobody jako pouhé nepřítomnosti zásahu (non-interference) pojetí náročnější. 
Svobodný nejsem tehdy, když do mého jednání nikdo právě nezasahuje, nýbrž tehdy, když nade mnou nikdo nedisponuje kapacitou svévolně zasáhnout.
 Rozdíl mezi oběma pojetími odhaluje klasický příklad, jenž stojí v jádru celé tradice: otrok laskavého pána. Pán, který svého otroka nikdy netrestá a ponechává mu široký prostor jednání, do jeho voleb fakticky nezasahuje — v liberálním smyslu je tedy otrok \uv{svobodný}.
  A přece cítíme, že svobodný není. Chybí mu nikoli nezasahování, nýbrž jistota: prostor, jejž požívá, není jeho právem, nýbrž pánovou libovůlí, kterou lze kdykoli odvolat.
   Otrok je nucen žít v ustavičné závislosti na cizí dobré vůli — a právě tato závislost, nikoli uskutečněný zásah, je oním bezprávím, jež republikánská tradice nazývá nadvládou.
Lovett tuto strukturu upřesňuje trojicí podmínek: nadvláda předpokládá (i) závislostní vztah, z něhož nelze bez značných nákladů vystoupit, (ii) nerovnováhu moci mezi stranami a (iii) možnost svévolného uplatnění této moci.
 Rozhodující je třetí podmínka a spolu s ní výklad pojmu svévole. 
 Zde se obě verze tradice mírně rozcházejí a je užitečné rozdíl přiznat.
  Podle procedurálního čtení je moc svévolná tehdy, není-li vázána sdílenými pravidly, jimž se podrobený může dovolat. 
  Podle Pettitova náročnějšího čtení je svévolná tehdy, není-li nucena sledovat zájmy a názory těch, jichž se týká — bez ohledu na to, zda se náhodou drží pravidel či nikoli.
   Pro naši analýzu je nosné právě Pettitovo pojetí, neboť činí zřejmým, proč benevolence nestačí: moc, která se chová vstřícně, avšak činí tak z vlastního rozhodnutí a nikoli z donucení respektovat zájmy podrobených, zůstává mocí nadvládnou.
Odtud plyne rozhodující obrat celé práce. Uživatel algoritmické platformy je vůči provozující korporaci v témže postavení jako otrok vůči laskavému pánovi.
 I když konkrétní aplikace algoritmu není namířena proti němu, ba i když je infrastruktura nastavena způsobem, jenž mu subjektivně prospívá, netěší se svobodě v republikánském smyslu. 
 Rozhoduje se totiž v prostředí, jehož parametry určuje nevolený soukromý subjekt disponující kapacitou toto prostředí kdykoli svévolně změnit — změnit, co uživatel uvidí, koho nalezne, čí hlas k němu dolehne — aniž by tato kapacita podléhala kontrole těch, kdo jsou jí vystaveni. 
 Nadvláda tedy nespočívá v uskutečněné škodě, nýbrž v samotné existenci takové nevázané, nekontrolované a demokraticky neodpovědné moci nad infrastrukturou, v níž se veřejná racionalita jedině může uskutečňovat.
Tímto vymezením získává práce tři výhody, jež strukturují celý následující výklad.
 

Za prvé, obvinění z nadvlády je robustní vůči empirické revizi: i kdyby se prokázalo, že algoritmická média nezpůsobují měřitelnou kognitivní škodu, bezpráví nadvlády by trvalo, neboť nezávisí na tom, zda a jak je moc uplatněna. 
 

Za druhé, přesouvá otázku z roviny terapeutické (jak chránit jednotlivce před závislostí a manipulací) do roviny ústavně-politické (jak moc rozdělit a učinit ji odpovědnou). 


Za třetí, vysvětluje strukturální selhání stávající regulace: individualisticky pojatá ochrana práv — souhlas, přístup k datům, právo na vysvětlení — nemůže nadvládu odstranit, neboť ta je vztahem, a nikoli událostí.
 Informovaný souhlas jednotlivce nadvládný vztah neruší, nýbrž jej stvrzuje; k jeho odstranění je třeba změnit rozložení moci samo. 
 K institucionálním důsledkům tohoto zjištění se práce vrací ve třetí kapitole.